Los mockups son representaciones visuales de la interfaz de usuario de un proyecto que 
anticipar y evaluar el diseño final antes de su implementación. Su objetivo es ofrecer una visión clara de los elementos gráficos, de la 
estructura y la disposición de los componentes de la aplicación. 

Para el desarrollo de los mockups, se ha utilizado la herramienta Figma. A continuación se detallan 
las principales vistas que se han realizado. 

\paragraph{Pantalla principal}
En la figura \ref{FIG:FIGMAA} se muestra la pantalla principal de la aplicación a la que accede el usuario tras autenticarse. 
En la parte superior se encuentra una barra de navegación que permite navegar entre las distintas
secciones de la aplicación. Justo debajo, se encuentra una barra de búsqueda acompañada un boton desplegable 
con distintos filtros como se observa figura \ref{FIG:FIGMA}. Por último, se presenta un listado de videojuegos mostrados 
mediante sus imágenes, con soporte para paginación.

\begin{figure}[Vista principal]{FIG:FIGMAA}{Vista principal de la aplicación}
    \image{10cm}{}{mac1}
\end{figure}

\begin{figure}[Vista principal con filtros]{FIG:FIGMA}{Vista principal de la aplicación con filtros desplegados}
    \image{10cm}{}{mac2}
\end{figure}

\paragraph{Opciones videojuegos}
En la figura \ref{FIG:GAME} se muestra la vista con la información detallada del videojuego, a la que accede al usuario tras hacer clic en 
su imagen. En esta pantalla se presentan distintas opciones como añadir el videojuego a la wishlist, marcarlo como jugado (played list) o 
realizar una valoración. Estas acciones están disponibles mediante iconos situados a la derecha de la imagen del videojuego. En orden descendente, 
las acciones asociadas a los iconos son: añadir o eliminar videojuegos de la wishlist, añadir o eliminar videojuegos de la played list y realizar una valoración.
A la izquierda de la imagen, cuenta con un botón para volver a la página anterior en la que se encontraba el usuario.

\begin{figure}[Información del videojuego]{FIG:GAME}{Vista donde se muestra la información del videojuego}
    \image{8cm}{}{game}
\end{figure}

\paragraph{Listas personalizadas}
En la figura \ref{FIG:LISTAS} se ilustra la vista correspondiente a las listas personalizadas del usuario. En esta pantalla, el usuario puede crear 
una nueva lista indicando un nombre y los videojuegos que quiere asociar, tal y como se muestra en la figura \ref{FIG:LISTACREAR}. Además, puede consultar 
el contenido de cada lista, así como eliminarla o editarla.

\begin{figure}[Listas personalizadas]{FIG:LISTAS}{Vista donde se muestran las listas personalizadas y las diferentes acciones disponibles como eliminar y editar}
    \image{8cm}{}{lista}
\end{figure}

\begin{figure}[Creación de listas personalizadas]{FIG:LISTACREAR}{Vista donde se muestra un formulario para la creación de vistas personalizadas}
    \image{8cm}{}{crear}
\end{figure}

\paragraph{Amistad}
En la figura \ref{FIG:AMISTAD} se muestra la pantalla correspondiente a las amistades del usuario. En esta vista se divide en dos secciones principales: 
una en la que se listan los usuarios con los que ha establecido amistad, y otra en la que se muestran las solicitudes de amistad pendientes, como 
se ilustra en la figura \ref{FIG:SOLICITUD}

\begin{figure}[Amistad entre usuarios]{FIG:AMISTAD}{Vista donde se muestra los usuarios con lo que tiene amistad}
    \image{10cm}{}{amistad}
\end{figure}

\begin{figure}[Solicitudes de amistad pendientes]{FIG:SOLICITUD}{Vista donde se muestran las solicitudes de amistad pendientes y las diferentes acciones disponibles como aceptar y eliminar }
    \image{10cm}{}{solicitud}
\end{figure}

\paragraph{Actividad}
Por último, la figura \ref{FIG:ACTIVIDAD} se muestra la pantalla donde se visualiza la actividad de otros usuarios. Como se observa, se incluye las 
valoraciones que ha realizado dicho usuario incluyendo la puntuación y el comentario, si existe.

\begin{figure}[Actividad de otros usuarios]{FIG:ACTIVIDAD}{Vista donde se muestra la actividad de otros usuarios }
    \image{10cm}{}{actividad}
\end{figure}